% paper24-async-effects.tex — Async Effect Boundaries: Verifying Python's
%   async/await Through Sheaf Decomposition
% Judgment Geometry Series, Paper 24 (async effects)
\documentclass[11pt]{article}
% jugeo-common.tex — shared preamble for all Judgment Geometry papers
% Include via: % jugeo-common.tex — shared preamble for all Judgment Geometry papers
% Include via: % jugeo-common.tex — shared preamble for all Judgment Geometry papers
% Include via: \input{jugeo-common}

% ── Packages ────────────────────────────────────────────────────────────
\usepackage{amsmath,amssymb,amsthm,mathtools}
\usepackage{tikz-cd}
\usepackage{listings}
\usepackage{xcolor}
\usepackage{xspace}
\usepackage{booktabs}
\usepackage{multirow}
\usepackage{enumitem}
\usepackage[expansion=false]{microtype}
\usepackage{hyperref}
\usepackage{cleveref}
% \usepackage{thmtools}  % disabled: not available in this TeX installation
\usepackage{float}
\usepackage{caption}
\usepackage{subcaption}
\usepackage{tabularx}
\usepackage{stmaryrd}       % \llbracket, \rrbracket
\usepackage{bbm}            % \mathbbm{1}

% ── Page geometry ───────────────────────────────────────────────────────
\usepackage[margin=1in]{geometry}

% ── Theorem environments ────────────────────────────────────────────────
\theoremstyle{plain}
\newtheorem{theorem}{Theorem}[section]
\newtheorem{lemma}[theorem]{Lemma}
\newtheorem{proposition}[theorem]{Proposition}
\newtheorem{corollary}[theorem]{Corollary}

\theoremstyle{definition}
\newtheorem{definition}[theorem]{Definition}
\newtheorem{example}[theorem]{Example}
\newtheorem{construction}[theorem]{Construction}

\theoremstyle{remark}
\newtheorem{remark}[theorem]{Remark}
\newtheorem{notation}[theorem]{Notation}

% ── Python listings style ──────────────────────────────────────────────
\definecolor{jugeoBlue}{HTML}{2563EB}
\definecolor{jugeoGreen}{HTML}{059669}
\definecolor{jugeoRed}{HTML}{DC2626}
\definecolor{jugeoPurple}{HTML}{7C3AED}
\definecolor{jugeoGray}{HTML}{6B7280}
\definecolor{jugeoBg}{HTML}{F8FAFC}

\lstdefinestyle{jugeo-python}{
  language=Python,
  basicstyle=\ttfamily\small,
  keywordstyle=\color{jugeoBlue}\bfseries,
  stringstyle=\color{jugeoGreen},
  commentstyle=\color{jugeoGray}\itshape,
  numberstyle=\tiny\color{jugeoGray},
  backgroundcolor=\color{jugeoBg},
  numbers=left,
  numbersep=6pt,
  frame=single,
  framerule=0.4pt,
  rulecolor=\color{jugeoGray!40},
  breaklines=true,
  breakatwhitespace=true,
  showstringspaces=false,
  tabsize=4,
  xleftmargin=1.5em,
  framexleftmargin=1.5em,
  morekeywords={dataclass,frozen,field,Enum,IntEnum,auto,Any,
                Mapping,Sequence,Optional,match,case},
  literate={→}{{$\to$}}1 {←}{{$\leftarrow$}}1
           {≤}{{$\leq$}}1 {≥}{{$\geq$}}1 {≠}{{$\neq$}}1
           {∈}{{$\in$}}1 {∀}{{$\forall$}}1 {∃}{{$\exists$}}1
           {⊕}{{$\oplus$}}1 {⊖}{{$\ominus$}}1,
  escapeinside={(*@}{@*)},
}
\lstset{style=jugeo-python}

\lstdefinestyle{jugeo-output}{
  basicstyle=\ttfamily\small,
  backgroundcolor=\color{jugeoBg},
  frame=single,
  framerule=0.4pt,
  rulecolor=\color{jugeoGray!40},
  breaklines=true,
  showstringspaces=false,
  xleftmargin=1.5em,
  framexleftmargin=0.5em,
  numbers=none,
}

% ── JuGeo-specific macros ──────────────────────────────────────────────

% Judgment 8-tuple
\newcommand{\J}{\mathcal{J}}
\newcommand{\Jud}[8]{(#1,\,#2,\,#3,\,#4,\,#5,\,#6,\,#7,\,#8)}
\newcommand{\JudShort}[1]{\J_{#1}}

% Components
\newcommand{\coord}{\mathsf{c}}            % coordinate
\newcommand{\prop}{\varphi}                 % proposition
\newcommand{\carrier}{\mathcal{A}}          % carrier type
\newcommand{\evid}{\mathcal{E}}             % evidence bundle
\newcommand{\oblig}{\mathcal{O}}            % obligations
\newcommand{\obstr}{\mathcal{B}}            % obstructions
\newcommand{\trust}{\mathcal{T}}            % trust annotation
\newcommand{\prov}{\Pi}                     % provenance

% Trust algebra
\newcommand{\Talg}{\mathfrak{T}}            % trust ordered algebra
\newcommand{\Eadm}{\mathcal{E}_{\mathrm{adm}}}
\newcommand{\tleq}{\preceq}                 % trust ordering
\newcommand{\tjoin}{\oplus}                 % conservative join
\newcommand{\tatten}{\ominus}               % attenuation
\newcommand{\tpromote}{\uparrow_\pi}        % promotion
\newcommand{\tdemote}{\downarrow_\chi}       % demotion

% Trust tiers
\newcommand{\Tcontra}{\ensuremath{\mathsf{CONTRADICTED}}}
\newcommand{\Tunver}{\ensuremath{\mathsf{UNVERIFIED}}}
\newcommand{\Tcopilot}{\ensuremath{\mathsf{COPILOT}}}
\newcommand{\Toracle}{\ensuremath{\mathsf{ORACLE}}}
\newcommand{\Truntime}{\ensuremath{\mathsf{RUNTIME}}}
\newcommand{\Tsolver}{\ensuremath{\mathsf{SOLVER}}}
\newcommand{\Tproof}{\ensuremath{\mathsf{PROOF}}}

% Site and geometry
\newcommand{\Site}{\mathcal{S}}             % semantic site
\newcommand{\Cov}{\mathrm{Cov}}            % covering family
\newcommand{\Ob}{\mathrm{Ob}}              % objects
\newcommand{\Mor}{\mathrm{Mor}}            % morphisms
\newcommand{\res}{\mathrm{res}}            % restriction
\newcommand{\Cech}{\check{C}}              % Čech complex
\newcommand{\Hcoh}[1]{H^{#1}}             % cohomology group

% Descent
\newcommand{\descent}{\mathrm{desc}}
\newcommand{\glue}{\mathrm{glue}}
\newcommand{\overlap}[2]{#1 \cap #2}

% Semantic moves
\newcommand{\VERIFY}{\textsc{Verify}}
\newcommand{\CONSTRUCT}{\textsc{Construct}}
\newcommand{\NORMALIZE}{\textsc{Normalize}}
\newcommand{\GLUE}{\textsc{Glue}}
\newcommand{\RESTRICT}{\textsc{Restrict}}
\newcommand{\TRANSPORT}{\textsc{Transport}}
\newcommand{\REPAIR}{\textsc{Repair}}
\newcommand{\PROMOTE}{\textsc{Promote}}
\newcommand{\CHALLENGE}{\textsc{Challenge}}
\newcommand{\ESCALATE}{\textsc{Escalate}}

% Fragments
\newcommand{\QFLIA}{\texttt{QF\_LIA}}
\newcommand{\QFLRA}{\texttt{QF\_LRA}}
\newcommand{\QFBV}{\texttt{QF\_BV}}
\newcommand{\QFUF}{\texttt{QF\_UF}}

% Misc
\newcommand{\Python}{\textsc{Python}}
\newcommand{\JuGeo}{\textsc{JuGeo}}
\newcommand{\LEAN}{\textsc{Lean}\,4}
\newcommand{\Fstar}{F$^{\star}$}
\newcommand{\Coq}{\textsc{Coq}}
\newcommand{\Dafny}{\textsc{Dafny}}

% Common shortcuts
\newcommand{\ie}{i.e.\@\xspace}
\newcommand{\eg}{e.g.\@\xspace}
\newcommand{\cf}{cf.\@\xspace}
\newcommand{\wrt}{w.r.t.\@\xspace}

% Evaluation shortcuts
\newcommand{\numcases}{300}
\newcommand{\numspec}{100}
\newcommand{\numeq}{100}
\newcommand{\numbug}{100}
\newcommand{\medtime}{6.5\,ms}
\newcommand{\totaltime}{1.949\,s}


% ── Packages ────────────────────────────────────────────────────────────
\usepackage{amsmath,amssymb,amsthm,mathtools}
\usepackage{tikz-cd}
\usepackage{listings}
\usepackage{xcolor}
\usepackage{xspace}
\usepackage{booktabs}
\usepackage{multirow}
\usepackage{enumitem}
\usepackage[expansion=false]{microtype}
\usepackage{hyperref}
\usepackage{cleveref}
% \usepackage{thmtools}  % disabled: not available in this TeX installation
\usepackage{float}
\usepackage{caption}
\usepackage{subcaption}
\usepackage{tabularx}
\usepackage{stmaryrd}       % \llbracket, \rrbracket
\usepackage{bbm}            % \mathbbm{1}

% ── Page geometry ───────────────────────────────────────────────────────
\usepackage[margin=1in]{geometry}

% ── Theorem environments ────────────────────────────────────────────────
\theoremstyle{plain}
\newtheorem{theorem}{Theorem}[section]
\newtheorem{lemma}[theorem]{Lemma}
\newtheorem{proposition}[theorem]{Proposition}
\newtheorem{corollary}[theorem]{Corollary}

\theoremstyle{definition}
\newtheorem{definition}[theorem]{Definition}
\newtheorem{example}[theorem]{Example}
\newtheorem{construction}[theorem]{Construction}

\theoremstyle{remark}
\newtheorem{remark}[theorem]{Remark}
\newtheorem{notation}[theorem]{Notation}

% ── Python listings style ──────────────────────────────────────────────
\definecolor{jugeoBlue}{HTML}{2563EB}
\definecolor{jugeoGreen}{HTML}{059669}
\definecolor{jugeoRed}{HTML}{DC2626}
\definecolor{jugeoPurple}{HTML}{7C3AED}
\definecolor{jugeoGray}{HTML}{6B7280}
\definecolor{jugeoBg}{HTML}{F8FAFC}

\lstdefinestyle{jugeo-python}{
  language=Python,
  basicstyle=\ttfamily\small,
  keywordstyle=\color{jugeoBlue}\bfseries,
  stringstyle=\color{jugeoGreen},
  commentstyle=\color{jugeoGray}\itshape,
  numberstyle=\tiny\color{jugeoGray},
  backgroundcolor=\color{jugeoBg},
  numbers=left,
  numbersep=6pt,
  frame=single,
  framerule=0.4pt,
  rulecolor=\color{jugeoGray!40},
  breaklines=true,
  breakatwhitespace=true,
  showstringspaces=false,
  tabsize=4,
  xleftmargin=1.5em,
  framexleftmargin=1.5em,
  morekeywords={dataclass,frozen,field,Enum,IntEnum,auto,Any,
                Mapping,Sequence,Optional,match,case},
  literate={→}{{$\to$}}1 {←}{{$\leftarrow$}}1
           {≤}{{$\leq$}}1 {≥}{{$\geq$}}1 {≠}{{$\neq$}}1
           {∈}{{$\in$}}1 {∀}{{$\forall$}}1 {∃}{{$\exists$}}1
           {⊕}{{$\oplus$}}1 {⊖}{{$\ominus$}}1,
  escapeinside={(*@}{@*)},
}
\lstset{style=jugeo-python}

\lstdefinestyle{jugeo-output}{
  basicstyle=\ttfamily\small,
  backgroundcolor=\color{jugeoBg},
  frame=single,
  framerule=0.4pt,
  rulecolor=\color{jugeoGray!40},
  breaklines=true,
  showstringspaces=false,
  xleftmargin=1.5em,
  framexleftmargin=0.5em,
  numbers=none,
}

% ── JuGeo-specific macros ──────────────────────────────────────────────

% Judgment 8-tuple
\newcommand{\J}{\mathcal{J}}
\newcommand{\Jud}[8]{(#1,\,#2,\,#3,\,#4,\,#5,\,#6,\,#7,\,#8)}
\newcommand{\JudShort}[1]{\J_{#1}}

% Components
\newcommand{\coord}{\mathsf{c}}            % coordinate
\newcommand{\prop}{\varphi}                 % proposition
\newcommand{\carrier}{\mathcal{A}}          % carrier type
\newcommand{\evid}{\mathcal{E}}             % evidence bundle
\newcommand{\oblig}{\mathcal{O}}            % obligations
\newcommand{\obstr}{\mathcal{B}}            % obstructions
\newcommand{\trust}{\mathcal{T}}            % trust annotation
\newcommand{\prov}{\Pi}                     % provenance

% Trust algebra
\newcommand{\Talg}{\mathfrak{T}}            % trust ordered algebra
\newcommand{\Eadm}{\mathcal{E}_{\mathrm{adm}}}
\newcommand{\tleq}{\preceq}                 % trust ordering
\newcommand{\tjoin}{\oplus}                 % conservative join
\newcommand{\tatten}{\ominus}               % attenuation
\newcommand{\tpromote}{\uparrow_\pi}        % promotion
\newcommand{\tdemote}{\downarrow_\chi}       % demotion

% Trust tiers
\newcommand{\Tcontra}{\ensuremath{\mathsf{CONTRADICTED}}}
\newcommand{\Tunver}{\ensuremath{\mathsf{UNVERIFIED}}}
\newcommand{\Tcopilot}{\ensuremath{\mathsf{COPILOT}}}
\newcommand{\Toracle}{\ensuremath{\mathsf{ORACLE}}}
\newcommand{\Truntime}{\ensuremath{\mathsf{RUNTIME}}}
\newcommand{\Tsolver}{\ensuremath{\mathsf{SOLVER}}}
\newcommand{\Tproof}{\ensuremath{\mathsf{PROOF}}}

% Site and geometry
\newcommand{\Site}{\mathcal{S}}             % semantic site
\newcommand{\Cov}{\mathrm{Cov}}            % covering family
\newcommand{\Ob}{\mathrm{Ob}}              % objects
\newcommand{\Mor}{\mathrm{Mor}}            % morphisms
\newcommand{\res}{\mathrm{res}}            % restriction
\newcommand{\Cech}{\check{C}}              % Čech complex
\newcommand{\Hcoh}[1]{H^{#1}}             % cohomology group

% Descent
\newcommand{\descent}{\mathrm{desc}}
\newcommand{\glue}{\mathrm{glue}}
\newcommand{\overlap}[2]{#1 \cap #2}

% Semantic moves
\newcommand{\VERIFY}{\textsc{Verify}}
\newcommand{\CONSTRUCT}{\textsc{Construct}}
\newcommand{\NORMALIZE}{\textsc{Normalize}}
\newcommand{\GLUE}{\textsc{Glue}}
\newcommand{\RESTRICT}{\textsc{Restrict}}
\newcommand{\TRANSPORT}{\textsc{Transport}}
\newcommand{\REPAIR}{\textsc{Repair}}
\newcommand{\PROMOTE}{\textsc{Promote}}
\newcommand{\CHALLENGE}{\textsc{Challenge}}
\newcommand{\ESCALATE}{\textsc{Escalate}}

% Fragments
\newcommand{\QFLIA}{\texttt{QF\_LIA}}
\newcommand{\QFLRA}{\texttt{QF\_LRA}}
\newcommand{\QFBV}{\texttt{QF\_BV}}
\newcommand{\QFUF}{\texttt{QF\_UF}}

% Misc
\newcommand{\Python}{\textsc{Python}}
\newcommand{\JuGeo}{\textsc{JuGeo}}
\newcommand{\LEAN}{\textsc{Lean}\,4}
\newcommand{\Fstar}{F$^{\star}$}
\newcommand{\Coq}{\textsc{Coq}}
\newcommand{\Dafny}{\textsc{Dafny}}

% Common shortcuts
\newcommand{\ie}{i.e.\@\xspace}
\newcommand{\eg}{e.g.\@\xspace}
\newcommand{\cf}{cf.\@\xspace}
\newcommand{\wrt}{w.r.t.\@\xspace}

% Evaluation shortcuts
\newcommand{\numcases}{300}
\newcommand{\numspec}{100}
\newcommand{\numeq}{100}
\newcommand{\numbug}{100}
\newcommand{\medtime}{6.5\,ms}
\newcommand{\totaltime}{1.949\,s}


% ── Packages ────────────────────────────────────────────────────────────
\usepackage{amsmath,amssymb,amsthm,mathtools}
\usepackage{tikz-cd}
\usepackage{listings}
\usepackage{xcolor}
\usepackage{xspace}
\usepackage{booktabs}
\usepackage{multirow}
\usepackage{enumitem}
\usepackage[expansion=false]{microtype}
\usepackage{hyperref}
\usepackage{cleveref}
% \usepackage{thmtools}  % disabled: not available in this TeX installation
\usepackage{float}
\usepackage{caption}
\usepackage{subcaption}
\usepackage{tabularx}
\usepackage{stmaryrd}       % \llbracket, \rrbracket
\usepackage{bbm}            % \mathbbm{1}

% ── Page geometry ───────────────────────────────────────────────────────
\usepackage[margin=1in]{geometry}

% ── Theorem environments ────────────────────────────────────────────────
\theoremstyle{plain}
\newtheorem{theorem}{Theorem}[section]
\newtheorem{lemma}[theorem]{Lemma}
\newtheorem{proposition}[theorem]{Proposition}
\newtheorem{corollary}[theorem]{Corollary}

\theoremstyle{definition}
\newtheorem{definition}[theorem]{Definition}
\newtheorem{example}[theorem]{Example}
\newtheorem{construction}[theorem]{Construction}

\theoremstyle{remark}
\newtheorem{remark}[theorem]{Remark}
\newtheorem{notation}[theorem]{Notation}

% ── Python listings style ──────────────────────────────────────────────
\definecolor{jugeoBlue}{HTML}{2563EB}
\definecolor{jugeoGreen}{HTML}{059669}
\definecolor{jugeoRed}{HTML}{DC2626}
\definecolor{jugeoPurple}{HTML}{7C3AED}
\definecolor{jugeoGray}{HTML}{6B7280}
\definecolor{jugeoBg}{HTML}{F8FAFC}

\lstdefinestyle{jugeo-python}{
  language=Python,
  basicstyle=\ttfamily\small,
  keywordstyle=\color{jugeoBlue}\bfseries,
  stringstyle=\color{jugeoGreen},
  commentstyle=\color{jugeoGray}\itshape,
  numberstyle=\tiny\color{jugeoGray},
  backgroundcolor=\color{jugeoBg},
  numbers=left,
  numbersep=6pt,
  frame=single,
  framerule=0.4pt,
  rulecolor=\color{jugeoGray!40},
  breaklines=true,
  breakatwhitespace=true,
  showstringspaces=false,
  tabsize=4,
  xleftmargin=1.5em,
  framexleftmargin=1.5em,
  morekeywords={dataclass,frozen,field,Enum,IntEnum,auto,Any,
                Mapping,Sequence,Optional,match,case},
  literate={→}{{$\to$}}1 {←}{{$\leftarrow$}}1
           {≤}{{$\leq$}}1 {≥}{{$\geq$}}1 {≠}{{$\neq$}}1
           {∈}{{$\in$}}1 {∀}{{$\forall$}}1 {∃}{{$\exists$}}1
           {⊕}{{$\oplus$}}1 {⊖}{{$\ominus$}}1,
  escapeinside={(*@}{@*)},
}
\lstset{style=jugeo-python}

\lstdefinestyle{jugeo-output}{
  basicstyle=\ttfamily\small,
  backgroundcolor=\color{jugeoBg},
  frame=single,
  framerule=0.4pt,
  rulecolor=\color{jugeoGray!40},
  breaklines=true,
  showstringspaces=false,
  xleftmargin=1.5em,
  framexleftmargin=0.5em,
  numbers=none,
}

% ── JuGeo-specific macros ──────────────────────────────────────────────

% Judgment 8-tuple
\newcommand{\J}{\mathcal{J}}
\newcommand{\Jud}[8]{(#1,\,#2,\,#3,\,#4,\,#5,\,#6,\,#7,\,#8)}
\newcommand{\JudShort}[1]{\J_{#1}}

% Components
\newcommand{\coord}{\mathsf{c}}            % coordinate
\newcommand{\prop}{\varphi}                 % proposition
\newcommand{\carrier}{\mathcal{A}}          % carrier type
\newcommand{\evid}{\mathcal{E}}             % evidence bundle
\newcommand{\oblig}{\mathcal{O}}            % obligations
\newcommand{\obstr}{\mathcal{B}}            % obstructions
\newcommand{\trust}{\mathcal{T}}            % trust annotation
\newcommand{\prov}{\Pi}                     % provenance

% Trust algebra
\newcommand{\Talg}{\mathfrak{T}}            % trust ordered algebra
\newcommand{\Eadm}{\mathcal{E}_{\mathrm{adm}}}
\newcommand{\tleq}{\preceq}                 % trust ordering
\newcommand{\tjoin}{\oplus}                 % conservative join
\newcommand{\tatten}{\ominus}               % attenuation
\newcommand{\tpromote}{\uparrow_\pi}        % promotion
\newcommand{\tdemote}{\downarrow_\chi}       % demotion

% Trust tiers
\newcommand{\Tcontra}{\ensuremath{\mathsf{CONTRADICTED}}}
\newcommand{\Tunver}{\ensuremath{\mathsf{UNVERIFIED}}}
\newcommand{\Tcopilot}{\ensuremath{\mathsf{COPILOT}}}
\newcommand{\Toracle}{\ensuremath{\mathsf{ORACLE}}}
\newcommand{\Truntime}{\ensuremath{\mathsf{RUNTIME}}}
\newcommand{\Tsolver}{\ensuremath{\mathsf{SOLVER}}}
\newcommand{\Tproof}{\ensuremath{\mathsf{PROOF}}}

% Site and geometry
\newcommand{\Site}{\mathcal{S}}             % semantic site
\newcommand{\Cov}{\mathrm{Cov}}            % covering family
\newcommand{\Ob}{\mathrm{Ob}}              % objects
\newcommand{\Mor}{\mathrm{Mor}}            % morphisms
\newcommand{\res}{\mathrm{res}}            % restriction
\newcommand{\Cech}{\check{C}}              % Čech complex
\newcommand{\Hcoh}[1]{H^{#1}}             % cohomology group

% Descent
\newcommand{\descent}{\mathrm{desc}}
\newcommand{\glue}{\mathrm{glue}}
\newcommand{\overlap}[2]{#1 \cap #2}

% Semantic moves
\newcommand{\VERIFY}{\textsc{Verify}}
\newcommand{\CONSTRUCT}{\textsc{Construct}}
\newcommand{\NORMALIZE}{\textsc{Normalize}}
\newcommand{\GLUE}{\textsc{Glue}}
\newcommand{\RESTRICT}{\textsc{Restrict}}
\newcommand{\TRANSPORT}{\textsc{Transport}}
\newcommand{\REPAIR}{\textsc{Repair}}
\newcommand{\PROMOTE}{\textsc{Promote}}
\newcommand{\CHALLENGE}{\textsc{Challenge}}
\newcommand{\ESCALATE}{\textsc{Escalate}}

% Fragments
\newcommand{\QFLIA}{\texttt{QF\_LIA}}
\newcommand{\QFLRA}{\texttt{QF\_LRA}}
\newcommand{\QFBV}{\texttt{QF\_BV}}
\newcommand{\QFUF}{\texttt{QF\_UF}}

% Misc
\newcommand{\Python}{\textsc{Python}}
\newcommand{\JuGeo}{\textsc{JuGeo}}
\newcommand{\LEAN}{\textsc{Lean}\,4}
\newcommand{\Fstar}{F$^{\star}$}
\newcommand{\Coq}{\textsc{Coq}}
\newcommand{\Dafny}{\textsc{Dafny}}

% Common shortcuts
\newcommand{\ie}{i.e.\@\xspace}
\newcommand{\eg}{e.g.\@\xspace}
\newcommand{\cf}{cf.\@\xspace}
\newcommand{\wrt}{w.r.t.\@\xspace}

% Evaluation shortcuts
\newcommand{\numcases}{300}
\newcommand{\numspec}{100}
\newcommand{\numeq}{100}
\newcommand{\numbug}{100}
\newcommand{\medtime}{6.5\,ms}
\newcommand{\totaltime}{1.949\,s}


% ── Paper-specific macros ─────────────────────────────────────────────
\newcommand{\Async}{\mathsf{Async}}          % async region label
\newcommand{\Await}{\mathsf{await}}          % await point
\newcommand{\AwaitMor}[2]{\rho_{#1 \to #2}} % await restriction morphism
\newcommand{\CB}{\mathcal{CB}}               % concurrency boundary
\newcommand{\asyncSite}{\Site_{\Async}}      % async sub-site
\newcommand{\asyncCoord}[1]{\coord_{\Async}^{#1}} % async coordinate
\newcommand{\asyncSheaf}{\mathcal{F}_{\Async}} % async effect sheaf
\newcommand{\Epure}{\mathsf{Pure}}           % pure effect class
\newcommand{\Eexn}{\mathsf{Exn}}             % exception effect class
\newcommand{\Estate}{\mathsf{State}}         % stateful effect class
\newcommand{\Eio}{\mathsf{IO}}               % IO effect class
\newcommand{\effectClass}[1]{\mathsf{Eff}(#1)} % effect class of region
\newcommand{\asyncRegion}[1]{\alpha_{#1}}    % async region #1
\newcommand{\syncRegion}[1]{\sigma_{#1}}     % sync (caller) region #1
\newcommand{\scopeEnv}{\Gamma_{\mathsf{sc}}} % scope environment
\newcommand{\stateEnv}{\Sigma}               % state environment
\newcommand{\asyncSafe}{\mathsf{AsyncSafe}}  % safety predicate
\newcommand{\raceFree}{\mathsf{RaceFree}}    % race-freedom predicate
\newcommand{\boundary}{\partial}             % boundary operator
\newcommand{\sharedState}[1]{\hat{s}_{#1}}  % shared state variable
\newcommand{\taskGraph}{\mathcal{G}_{\mathsf{task}}} % task dependency graph
\newcommand{\cancelObs}{\mathcal{B}_{\mathsf{cancel}}} % cancellation obstruction
\newcommand{\asyncCov}{\mathrm{Cov}_{\Async}} % async covering
\newcommand{\gluedSec}{\tilde{s}}            % glued global section
\newcommand{\restriction}[2]{\res_{#1}^{#2}} % restriction map from #1 to #2

\title{\textbf{Async Effect Boundaries:\\
       Verifying Python's \texttt{async}/\texttt{await}\\
       Through Sheaf Decomposition}}
\author{%
  Halley Young\\
  \textit{Microsoft Research}\\
  \texttt{halleyyoung@microsoft.com}
}
\date{\today}

\begin{document}

\maketitle

\begin{abstract}
Python's \texttt{async}/\texttt{await} model introduces a form of cooperative
concurrency in which coroutines voluntarily yield control at \texttt{await}
points.  Verifying that an async program is race-free and effect-safe is
difficult because the interleaving structure is implicit in the control flow,
shared state may be mutated between suspension points, and cancellation
can inject exceptions at arbitrary positions in an otherwise pure coroutine.
We show that these challenges dissolve when each \emph{async region}---a
maximal block of code between two consecutive \texttt{await} points---is
treated as a coordinate on the \emph{async sub-site} $\asyncSite$, and each
\texttt{await} expression is treated as a \emph{restriction morphism} in the
Grothendieck topology.  The resulting \emph{async effect sheaf} $\asyncSheaf$
records four effect classes (pure, exception, stateful, IO) for every region;
\emph{concurrency boundaries} $\CB$ are the covers through which shared state
may escape across regions.  We prove the central theorem:
\emph{if every async region is locally verified, and every concurrency boundary
is clean, then the global program is race-free}.
The framework is implemented in the \JuGeo{} \Python{} runtime as the
\texttt{AsyncAnalyzer} and \texttt{ConcurrencyBoundary} components, and
integration with \texttt{ScopeAnalyzer} / \texttt{StateTracker} provides
ambient scope and state context for each region.
All main results are formalized in \LEAN{} (no \texttt{sorry}).
\end{abstract}

\tableofcontents
\bigskip

% ======================================================================
\section{Introduction}\label{sec:intro}
% ======================================================================

Python's coroutine machinery, introduced in PEP 492 and extended through
PEP 525 and PEP 567, provides a cooperative concurrency model in which
coroutines execute uninterrupted between successive \texttt{await}
expressions.  In contrast to preemptive threading, control transfers are
explicit: a running coroutine hands control back to the event loop only
when it hits an \texttt{await} site, making the interleaving structure
entirely deterministic given a fixed task schedule.

Yet async Python programs are notoriously difficult to verify.  Three sources
of complexity dominate:
\begin{enumerate}[label=(\roman*)]
  \item \textbf{Implicit shared state.}  A mutable object reachable from
        multiple coroutines may be read and written in any interleaving
        order, even though the event loop is single-threaded.  Race
        conditions arise not from data races in the traditional sense but
        from \emph{logical races}: a coroutine reads a value at one
        \texttt{await} point and acts on a stale view of it at the next.

  \item \textbf{Cancellation.}  Python's \texttt{asyncio.Task.cancel()} injects
        a \texttt{CancelledError} at an arbitrary \texttt{await} site.  A
        cancellation-safe coroutine must preserve its invariants even when
        interrupted mid-operation, a property that is hard to state let alone
        verify.

  \item \textbf{Effect heterogeneity.}  Async functions mix pure computation,
        IO operations, state mutations, and exception handling in a single
        control-flow graph.  Classical monadic effect systems handle one
        effect at a time; Python's coroutines routinely stack all four.
\end{enumerate}

The Judgment Geometry framework addresses these challenges through its
semantic-site machinery.  In Papers 1--3, we showed that code regions
become coordinates on a Grothendieck site $\Site$, and that properties
verified locally can be \emph{glued} to global results whenever appropriate
covering conditions hold.  The present paper applies this principle
specifically to async code.

\paragraph{Main contributions.}
\begin{enumerate}
  \item We define the \emph{async sub-site} $\asyncSite$ whose objects are
        \emph{async regions}---maximal straight-line code blocks between
        consecutive \texttt{await} points---and whose morphisms are
        \emph{await restriction maps} induced by the suspension points
        (\cref{sec:async-decomp,sec:await-morphisms}).

  \item We define the \emph{async effect sheaf} $\asyncSheaf$ and show it
        satisfies the sheaf axioms over the Grothendieck topology generated by
        the async covering families (\cref{sec:effect-class}).

  \item We formalize \emph{concurrency boundaries} as the covers through which
        shared mutable state crosses async region boundaries
        (\cref{sec:conc-boundaries}).

  \item We integrate scope and state tracking via \texttt{ScopeAnalyzer} and
        \texttt{StateTracker}, providing ambient lexical environments for each
        region (\cref{sec:scope-state}).

  \item We prove the \emph{Async Safety Theorem}: local verification of each
        region plus clean concurrency boundaries implies global race-freedom
        (\cref{sec:theorem}).
\end{enumerate}

\paragraph{Implementation.}
The ideas in this paper are realized in the \JuGeo{} Python runtime at
\texttt{src/jugeo/python\_runtime/effects\_async/} (the \texttt{AsyncAnalyzer}
component and related modules) and
\texttt{src/jugeo/python\_runtime/concurrency\_boundaries/} (the
\texttt{ConcurrencyBoundary} component).

\paragraph{Organization.}
\cref{sec:async-decomp} defines the async site;
\cref{sec:await-morphisms} treats await points as restriction morphisms;
\cref{sec:conc-boundaries} introduces concurrency boundaries;
\cref{sec:effect-class} develops the four-class effect system;
\cref{sec:scope-state} covers scope and state tracking;
\cref{sec:theorem} states and proves the Async Safety Theorem;
\cref{sec:experiments} reports experiments;
\cref{sec:conclusion} concludes.

% ======================================================================
\section{Async Decomposition}\label{sec:async-decomp}
% ======================================================================

\subsection{Async regions as coordinates}

Let $P$ be a Python coroutine defined with \texttt{async def}.  We say a
contiguous block of code $B \subseteq P$ is an \emph{async region} if:
\begin{enumerate}[label=(\alph*)]
  \item The first statement of $B$ is either the function entry or
        immediately follows an \texttt{await} expression.
  \item The last statement of $B$ either is the function exit or immediately
        precedes an \texttt{await} expression.
  \item No \texttt{await} expression appears in the interior of $B$.
\end{enumerate}
Condition (c) means $B$ executes without suspending the coroutine: from the
event loop's perspective, $B$ is \emph{atomic}.

\begin{definition}[Async coordinate]\label{def:async-coord}
Given a coroutine $P$, its set of async regions $\{\asyncRegion{1}, \dots,
\asyncRegion{n}\}$ gives rise to \emph{async coordinates}
$\asyncCoord{1}, \dots, \asyncCoord{n}$ in the semantic site, one per region.
Each coordinate is assigned the kind \texttt{region} in the
\texttt{CoordinateKind} vocabulary.
\end{definition}

\begin{example}\label{ex:simple-async}
Consider the coroutine:
\begin{lstlisting}[style=jugeo-python]
async def fetch_and_process(url: str, cache: dict) -> str:
    # Region alpha_1: entry to first await
    key = url.strip()
    if key in cache:
        return cache[key]

    # await point a_1
    raw = await http_get(url)          # (*)

    # Region alpha_2: between first and second await
    data = parse(raw)
    cache[key] = data

    # await point a_2
    result = await validate(data)      # (**)

    # Region alpha_3: second await to exit
    return result
\end{lstlisting}
This coroutine has three async regions:
$\asyncRegion{1}$ (lines 2--5),
$\asyncRegion{2}$ (lines 9--11),
$\asyncRegion{3}$ (line 14).
The await points $a_1$ and $a_2$ are the transitions.
\end{example}

\subsection{The async sub-site}

\begin{definition}[Async sub-site]\label{def:async-site}
The \emph{async sub-site} $\asyncSite = (\mathcal{C}_{\Async}, J_{\Async})$
consists of:
\begin{itemize}
  \item \textbf{Objects.}  All async coordinates $\asyncCoord{i}$ across all
        coroutines in the program.
  \item \textbf{Morphisms.}  The \emph{await restriction morphisms}
        $\AwaitMor{i}{j}$ described in \cref{sec:await-morphisms}, plus
        identity morphisms.
  \item \textbf{Grothendieck topology $J_{\Async}$.}  A sieve $S$ on
        $\asyncCoord{i}$ is a covering sieve if the union of the source regions
        of all morphisms in $S$ covers $\asyncCoord{i}$ in the program-flow
        sense: every predecessor async region contributes to $S$.
\end{itemize}
\end{definition}

\begin{remark}
The topology $J_{\Async}$ is coarser than the topology on the full site
$\Site$ from Paper~1.  Covering families only require predecessor coverage,
not the full Čech coverage used for specifications.  This reflects the
causal nature of coroutine execution: a region's state is determined by its
immediate predecessors, not by arbitrary subsets of the site.
\end{remark}

\subsection{Analysis by AsyncAnalyzer}

The \texttt{AsyncAnalyzer} component traverses the AST of each coroutine
and constructs the async sub-site automatically.  It records, for each
region:
\begin{itemize}
  \item The list of \texttt{await} expressions entering and exiting the region.
  \item All free variables read or written in the region.
  \item The exception handlers whose scope includes the region.
  \item Whether the region is at the top level of the coroutine or nested
        inside a loop or conditional.
\end{itemize}

\begin{construction}[AsyncAnalyzer output]\label{con:async-analyzer}
Running \texttt{AsyncAnalyzer} on a coroutine $P$ produces:
\begin{align*}
  \texttt{regions}(P) &= \{\asyncRegion{1}, \dots, \asyncRegion{n}\}\\
  \texttt{awaits}(P)  &= \{a_1, \dots, a_{n-1}\}\\
  \texttt{freeVars}(\asyncRegion{i}) &\subseteq \textsc{Var}
\end{align*}
together with the flow graph $\asyncRegion{i} \xrightarrow{a_i} \asyncRegion{i+1}$
for each consecutive pair.
\end{construction}

% ======================================================================
\section{Await Morphisms}\label{sec:await-morphisms}
% ======================================================================

\subsection{Await points as restriction morphisms}

In the sheaf-theoretic picture, a \emph{restriction morphism}
$\res_U^V : \mathcal{F}(U) \to \mathcal{F}(V)$ transports sections from a
larger open $U$ to a smaller open $V \subseteq U$.  In our setting, an
\texttt{await} expression plays exactly this role: it transfers the
\emph{observable state} of the coroutine from the current region to the
successor region, discarding anything that does not survive the suspension.

\begin{definition}[Await restriction morphism]\label{def:await-mor}
Let $\asyncRegion{i}$ and $\asyncRegion{i+1}$ be consecutive async regions
separated by the await point $a_i$.  The \emph{await restriction morphism}
\[
  \AwaitMor{i}{i+1} : \asyncCoord{i} \longrightarrow \asyncCoord{i+1}
\]
is a morphism of kind \texttt{restriction} with:
\begin{itemize}
  \item \textbf{Source}: $\asyncCoord{i}$ (the region producing the awaitable).
  \item \textbf{Target}: $\asyncCoord{i+1}$ (the region consuming the result).
  \item \textbf{Data carried}: the return value of the awaitable plus any
        captured free variables that are still live after the suspension.
\end{itemize}
\end{definition}

\subsection{Functoriality}

Restriction morphisms must compose: if $\asyncRegion{i}$ flows to
$\asyncRegion{j}$ and $\asyncRegion{j}$ flows to $\asyncRegion{k}$, the
composition $\AwaitMor{j}{k} \circ \AwaitMor{i}{j}$ is the restriction
from $\asyncRegion{i}$ to $\asyncRegion{k}$ skipping the intermediate region.
In Python, this corresponds to a chain of consecutive \texttt{await}
expressions.

\begin{proposition}[Await morphisms form a category]\label{prop:await-cat}
The async coordinates and await restriction morphisms form a small category
$\mathcal{C}_{\Async}$.  Composition is associative and each coordinate
has an identity morphism.
\end{proposition}
\begin{proof}
Composition is given by transitivity of the flow relation on async regions.
Associativity follows from the associativity of path composition in the
control-flow graph.  Identity morphisms exist by the reflexivity of the
flow relation (an empty \texttt{await} chain from a region to itself).
\end{proof}

\subsection{Compatibility conditions at await points}

A section $s \in \asyncSheaf(\asyncCoord{i})$ is \emph{compatible} with
$s' \in \asyncSheaf(\asyncCoord{i+1})$ at the await point $a_i$ if:
\[
  \AwaitMor{i}{i+1}(s) = s'.
\]
This is the standard sheaf locality condition: the restriction of the
section on the larger region must agree with the section on the smaller
region at their common interface.

\begin{example}[Continuation of \cref{ex:simple-async}]
In \cref{ex:simple-async}, the variable \texttt{cache} is live in all three
regions.  The section $s_1 \in \asyncSheaf(\asyncCoord{1})$ records that
\texttt{cache} is read-only in $\asyncRegion{1}$.  After the await $a_1$,
the section $s_2 \in \asyncSheaf(\asyncCoord{2})$ records that \texttt{cache}
is written in $\asyncRegion{2}$.  The compatibility condition at $a_1$
asserts that $s_2$ is consistent with $s_1$: specifically, the write to
\texttt{cache} in $\asyncRegion{2}$ is the only way the variable changes.
\end{example}

% ======================================================================
\section{Concurrency Boundaries}\label{sec:conc-boundaries}
% ======================================================================

\subsection{Definition}

While each async region is atomic with respect to the event loop, two
coroutines running in the same thread share all objects in the heap.  A
\emph{concurrency boundary} is a point at which the heap contents visible
to one coroutine may differ from what another coroutine last wrote, because
the event loop interleaved their executions.

\begin{definition}[Concurrency boundary]\label{def:cb}
A \emph{concurrency boundary} $\CB = (V, \asyncRegion{i}, \asyncRegion{j})$
consists of:
\begin{itemize}
  \item A set $V \subseteq \textsc{Var}$ of \emph{shared variables}.
  \item Two async regions $\asyncRegion{i}$, $\asyncRegion{j}$ (not
        necessarily in the same coroutine) such that both regions read or
        write some $v \in V$.
  \item The property that $\asyncRegion{i}$ and $\asyncRegion{j}$ may be
        interleaved by the event loop: there exists a schedule in which
        $\asyncRegion{i}$ writes $v$ and $\asyncRegion{j}$ reads $v$ after
        the write, or vice versa.
\end{itemize}
\end{definition}

\begin{definition}[Clean concurrency boundary]\label{def:clean-cb}
A concurrency boundary $\CB = (V, \asyncRegion{i}, \asyncRegion{j})$ is
\emph{clean} if for every shared variable $v \in V$:
\begin{enumerate}[label=(\roman*)]
  \item At least one of $\asyncRegion{i}$, $\asyncRegion{j}$ treats $v$ as
        read-only; or
  \item Both regions treat $v$ as write-only (no read after write); or
  \item The variable $v$ is protected by a mutual-exclusion mechanism
        (\eg an \texttt{asyncio.Lock}) that is held across both regions.
\end{enumerate}
\end{definition}

\subsection{Boundaries as cover morphisms}

In sheaf language, a concurrency boundary corresponds to a \emph{covering
family}: $\{\asyncCoord{i}, \asyncCoord{j}\} \to \asyncCoord{v}$ where
$\asyncCoord{v}$ is the ``shared variable coordinate''---the coordinate
representing the heap location of $v$.  For the covering family to define
a clean cover (one that allows section gluing), the sections
$s_i \in \asyncSheaf(\asyncCoord{i})$ and $s_j \in \asyncSheaf(\asyncCoord{j})$
must agree on the overlap, \ie on the shared variable coordinate.

\begin{proposition}[Boundary and section compatibility]\label{prop:boundary-compat}
A concurrency boundary $\CB = (V, \asyncRegion{i}, \asyncRegion{j})$ is clean
if and only if the sections
$\restriction{\asyncCoord{i}}{\asyncCoord{v}}(s_i)$
and
$\restriction{\asyncCoord{j}}{\asyncCoord{v}}(s_j)$
agree for every $v \in V$.
\end{proposition}
\begin{proof}
($\Rightarrow$) If $\CB$ is clean by condition (i), $v$ is read-only in
one region, so its section records no write; the restriction to
$\asyncCoord{v}$ carries only the read type, which agrees with the
read-only section of the other region.  Conditions (ii) and (iii) follow
similarly by the definition of section agreement.
($\Leftarrow$) If the restricted sections agree, the only consistent
interpretations are those listed in \cref{def:clean-cb}, by the effect
classification of \cref{sec:effect-class}.
\end{proof}

\subsection{ConcurrencyBoundary in the implementation}

The \texttt{ConcurrencyBoundary} component, located at
\texttt{concurrency\_boundaries/}, identifies all pairs of async regions
across all simultaneously-active coroutines and tests each pair for the
clean-boundary conditions.  It outputs:
\begin{itemize}
  \item A list of all concurrency boundaries together with the shared
        variable set $V$ for each.
  \item A boolean \texttt{is\_clean} flag for each boundary.
  \item For dirty boundaries, a diagnostic message identifying which
        variable and which pair of regions violate the condition.
\end{itemize}

\begin{example}[Dirty boundary]\label{ex:dirty}
\begin{lstlisting}[style=jugeo-python]
counter: int = 0

async def increment():
    global counter
    # Region alpha_1
    tmp = counter        # read
    # await a_1
    await asyncio.sleep(0)
    # Region alpha_2
    counter = tmp + 1    # write-back (stale!)

async def main():
    await asyncio.gather(increment(), increment())
\end{lstlisting}
Here \texttt{counter} is shared between two concurrent \texttt{increment()}
tasks.  The boundary $({\{\texttt{counter}\}}, \asyncRegion{1}^{(1)},
\asyncRegion{1}^{(2)})$ is \emph{dirty}: both tasks read and write
\texttt{counter} without a lock.  The \texttt{ConcurrencyBoundary} component
flags this as a potential logical race.
\end{example}

% ======================================================================
\section{Effect Classification}\label{sec:effect-class}
% ======================================================================

\subsection{The four effect classes}

The Judgment Geometry effect system for async code uses four classes:

\begin{definition}[Effect classes]\label{def:effect-classes}
\begin{align*}
  \Epure   &: \text{no side effects; return value determined by arguments}\\
  \Eexn    &: \text{may raise exceptions (including \texttt{CancelledError})}\\
  \Estate  &: \text{reads or writes shared mutable state}\\
  \Eio     &: \text{performs IO (network, file system, signals)}
\end{align*}
These classes are ordered by subsumption:
$\Epure \sqsubseteq \Eexn \sqsubseteq \Estate \sqsubseteq \Eio$.
\end{definition}

\begin{remark}
The ordering reflects the degree of \emph{trust attenuation}: a region
classified as $\Estate$ may also raise exceptions, so $\Eexn \sqsubseteq
\Estate$.  Similarly, IO inherently involves state (file descriptors,
socket buffers) so $\Estate \sqsubseteq \Eio$.  This ordering mirrors the
trust lattice of the \JuGeo{} trust algebra (Paper~4).
\end{remark}

\subsection{The async effect sheaf}

\begin{definition}[Async effect sheaf]\label{def:async-sheaf}
The \emph{async effect sheaf} $\asyncSheaf$ on $\asyncSite$ assigns to each
async coordinate $\asyncCoord{i}$ the set:
\[
  \asyncSheaf(\asyncCoord{i}) = \{\Epure, \Eexn, \Estate, \Eio\}
\]
with restriction maps defined by:
$\AwaitMor{i}{i+1}^* : \asyncSheaf(\asyncCoord{i}) \to \asyncSheaf(\asyncCoord{i+1})$
given by $e \mapsto e \sqcup \effectClass{\asyncRegion{i+1}}$
(the effect of the successor region is at least as large as the propagated
effect class).
\end{definition}

\begin{proposition}[Sheaf axioms]\label{prop:sheaf-axioms}
$\asyncSheaf$ satisfies the sheaf axioms for the topology $J_{\Async}$:
\begin{enumerate}[label=(\roman*)]
  \item \textbf{Locality}: if two sections $s, s'$ on $\asyncCoord{i}$
        restrict to the same section on every member of a covering sieve,
        then $s = s'$.
  \item \textbf{Gluing}: given a compatible family of sections on a
        covering sieve, there exists a unique global section extending them.
\end{enumerate}
\end{proposition}
\begin{proof}
(i) The effect class is determined by the effects performed in the
region $\asyncRegion{i}$ itself; two sections that agree on all predecessor
restrictions must have the same root classification.
(ii) Effect classes form a join-semilattice; given a compatible family
$\{e_U\}$ on a covering sieve, define $\gluedSec = \bigsqcup_U e_U$.
This is the unique minimal upper bound, which is the unique glued section.
\end{proof}

\subsection{Annotation propagation}

Effect annotations propagate forward through await morphisms:
\[
  \effectClass{\asyncRegion{i+1}} \;=\;
  \effectClass{\asyncRegion{i}} \sqcup \effectClass{\text{awaitable}_{a_i}}
    \sqcup \effectClass{\asyncRegion{i+1}^{\text{local}}}
\]
where $\effectClass{\text{awaitable}_{a_i}}$ is the effect class of the
awaitable expression at the \texttt{await} point $a_i$, and
$\effectClass{\asyncRegion{i+1}^{\text{local}}}$ is the locally-inferred
effect class of region $\asyncRegion{i+1}$ ignoring upstream effects.

The \texttt{AsyncAnalyzer} performs this propagation as a forward dataflow
analysis over the async flow graph, producing a fully annotated version of
the async sub-site $\asyncSite$.

\subsection{Handling cancellation}

Cancellation is modeled as an $\Eexn$ injection at every \texttt{await}
point.  A region $\asyncRegion{i}$ is \emph{cancellation-safe} if its
local effect class satisfies:
\[
  \effectClass{\asyncRegion{i}} \in \{\Epure, \Eexn\}
  \;\text{ or }\;
  \asyncRegion{i} \text{ is guarded by a \texttt{try}/\texttt{finally} block}.
\]
The \texttt{CancelledError} is an obstruction class in the Čech cohomology of
the task graph $\taskGraph$: it prevents the gluing of sections across the
cancellation boundary.

\begin{definition}[Cancellation obstruction]\label{def:cancel-obs}
The \emph{cancellation obstruction} $\cancelObs$ is the Čech 1-cocycle
that assigns to each await morphism $\AwaitMor{i}{i+1}$ the boolean value
$1$ (obstruction present) if the awaitable at $a_i$ may raise
\texttt{CancelledError}, and $0$ otherwise.
\end{definition}

% ======================================================================
\section{Scope and State Tracking}\label{sec:scope-state}
% ======================================================================

\subsection{The scope environment}

The \texttt{ScopeAnalyzer} component (\texttt{scope\_and\_state/}) augments
the async sub-site with lexical scope information.  For each async region
$\asyncRegion{i}$, it computes:
\begin{align*}
  \texttt{locals}(\asyncRegion{i})  &\subseteq \textsc{Var}
    \quad (\text{variables assigned in } \asyncRegion{i})\\
  \texttt{frees}(\asyncRegion{i})   &\subseteq \textsc{Var}
    \quad (\text{free variables used in } \asyncRegion{i})\\
  \texttt{globals}(\asyncRegion{i}) &\subseteq \textsc{Var}
    \quad (\text{global names accessed in } \asyncRegion{i})
\end{align*}
These sets constitute the \emph{scope environment}
$\scopeEnv(\asyncRegion{i})$.

\begin{definition}[Scope coordinate]\label{def:scope-coord}
Each async region $\asyncRegion{i}$ is assigned a \emph{scope coordinate}
$\coord_{\mathsf{sc}}^i$ that extends the async coordinate with the scope
information:
\[
  \coord_{\mathsf{sc}}^i
  \;=\;
  \bigl(\asyncCoord{i},\; \scopeEnv(\asyncRegion{i})\bigr).
\]
\end{definition}

\subsection{The state environment}

The \texttt{StateTracker} component records, at each \texttt{await} point
$a_k$, the \emph{state environment} $\stateEnv(a_k)$: a map from shared
variable names to their possible types and values as inferred by
flow-sensitive analysis.  Formally:
\[
  \stateEnv(a_k) : \textsc{Var} \rightharpoonup \textsc{Type} \times \textsc{Values}
\]
where $\rightharpoonup$ denotes a partial function (not all variables need be
tracked).

\begin{definition}[State section]\label{def:state-sec}
The \emph{state section} at region $\asyncRegion{i}$ is the pair
$(s_{\mathrm{in}}^i, s_{\mathrm{out}}^i)$ where:
\begin{itemize}
  \item $s_{\mathrm{in}}^i = \stateEnv(a_{i-1})$ (state entering the region,
        from the previous await point; $\stateEnv(\bot)$ at the entry region).
  \item $s_{\mathrm{out}}^i = \stateEnv(a_i)$ (state exiting the region,
        at the next await point; $\stateEnv(\top)$ at the exit region).
\end{itemize}
\end{definition}

\subsection{Compatibility with the async sheaf}

The state environment must be compatible with the effect classification:
\begin{itemize}
  \item If $\effectClass{\asyncRegion{i}} = \Epure$, then
        $s_{\mathrm{in}}^i = s_{\mathrm{out}}^i$
        (no state change in a pure region).
  \item If $\effectClass{\asyncRegion{i}} = \Eexn$, then the only state
        changes are those induced by unwinding the stack during exception
        handling (local variables go out of scope).
  \item If $\effectClass{\asyncRegion{i}} = \Estate$ or $\Eio$, then
        $s_{\mathrm{out}}^i$ may differ from $s_{\mathrm{in}}^i$ on any
        variable in $\texttt{globals}(\asyncRegion{i})$.
\end{itemize}

\begin{proposition}[State-effect compatibility]\label{prop:state-eff}
Let $\effectClass{\asyncRegion{i}} = \Epure$.  Then for every shared
variable $v \in \texttt{globals}(\asyncRegion{i})$,
$s_{\mathrm{in}}^i(v) = s_{\mathrm{out}}^i(v)$.
\end{proposition}
\begin{proof}
A pure region performs no side effects; in particular it writes no shared
state.  Thus the \texttt{StateTracker} records no state change for~$v$
through $\asyncRegion{i}$.
\end{proof}

\subsection{Async context variables}

Python's \texttt{contextvars} module provides \emph{task-local} context
variables: each coroutine inherits a copy of the context at creation time,
and writes to context variables in one coroutine do not affect others.
Context variables are therefore \emph{locally pure} at concurrency
boundaries: a concurrency boundary on a context variable is always clean
by condition (i) of \cref{def:clean-cb} (the write in one coroutine is
invisible to the other).

The \texttt{ScopeAnalyzer} distinguishes context variables from ordinary
global variables and marks their boundaries as clean automatically,
reducing the number of dirty boundaries to report.

% ======================================================================
\section{The Async Safety Theorem}\label{sec:theorem}
% ======================================================================

\subsection{Statement}

We can now state the central result of this paper.

\begin{definition}[Locally verified region]\label{def:locally-verified}
An async region $\asyncRegion{i}$ is \emph{locally verified} if:
\begin{enumerate}[label=(\roman*)]
  \item Its effect class $\effectClass{\asyncRegion{i}}$ is correctly
        inferred by \texttt{AsyncAnalyzer}.
  \item Its state section $(s_{\mathrm{in}}^i, s_{\mathrm{out}}^i)$ is
        consistent with $\effectClass{\asyncRegion{i}}$ in the sense of
        \cref{prop:state-eff}.
  \item If $\effectClass{\asyncRegion{i}} \in \{\Eexn, \Estate, \Eio\}$, all
        exceptions that may propagate out of the region are listed in the
        region's \emph{obligation set} $\oblig(\asyncRegion{i})$.
\end{enumerate}
\end{definition}

\begin{definition}[Race-free program]\label{def:race-free}
An async program $P$ (a set of coroutines) is \emph{race-free} if there
exists no schedule of the event loop such that two coroutines simultaneously
perform conflicting operations (a read and a write, or two writes) on the
same shared variable without synchronization.
\end{definition}

\begin{theorem}[Async Safety Theorem]\label{thm:async-safety}
Let $P$ be an async program whose async sub-site $\asyncSite$ has been
constructed by \texttt{AsyncAnalyzer}.  Suppose:
\begin{enumerate}[label=(\roman*)]
  \item Every async region $\asyncRegion{i}$ in every coroutine of $P$
        is locally verified.
  \item Every concurrency boundary identified by \texttt{ConcurrencyBoundary}
        is clean.
\end{enumerate}
Then $P$ is race-free.
\end{theorem}
\begin{proof}
We proceed by showing that the global state section $\gluedSec$ exists and
is unique, which implies no conflicting operations are possible.

\textbf{Step 1: Local sections exist.}  Since each region $\asyncRegion{i}$
is locally verified, \cref{def:locally-verified}(ii) provides a compatible
pair $(s_{\mathrm{in}}^i, s_{\mathrm{out}}^i)$.  This is the local section
$s_i \in \asyncSheaf(\asyncCoord{i})$.

\textbf{Step 2: Compatibility at await morphisms.}  Fix two consecutive
regions $\asyncRegion{i}, \asyncRegion{i+1}$ in the same coroutine.
The restriction map $\AwaitMor{i}{i+1}^*$ sends $s_i$ to a section on
$\asyncCoord{i+1}$; by the definition of the state section, this is
exactly $s_{\mathrm{in}}^{i+1} = \stateEnv(a_i) = s_{\mathrm{out}}^i$.
So $\AwaitMor{i}{i+1}^*(s_i) = s_{i+1}|_{\asyncCoord{i}}$, establishing
compatibility.

\textbf{Step 3: Compatibility at concurrency boundaries.}  Fix a concurrency
boundary $\CB = (V, \asyncRegion{i}, \asyncRegion{j})$ (possibly from
different coroutines).  By hypothesis (ii), $\CB$ is clean.  By
\cref{prop:boundary-compat}, the restrictions of $s_i$ and $s_j$ to
$\asyncCoord{v}$ agree for every $v \in V$.

\textbf{Step 4: Gluing.}  The sections $\{s_i\}$ form a compatible family
over the covering $\asyncCov$ of the full async sub-site.  By
\cref{prop:sheaf-axioms}(ii), there exists a unique global section
$\gluedSec \in \asyncSheaf(\asyncSite)$ restricting to each $s_i$.

\textbf{Step 5: Race-freedom.}  The global section $\gluedSec$ assigns a
consistent effect class and state value to every shared variable across all
coroutines.  In particular, no two regions hold conflicting write-write or
read-write sections on the same variable, since such a conflict would violate
the agreement condition in Step~3 and prevent gluing.  Hence $P$ is
race-free.
\end{proof}

\subsection{Corollaries}

\begin{corollary}[Modular verification]\label{cor:modular}
To verify that $P$ is race-free, it suffices to:
\begin{enumerate}
  \item Verify each async region independently (local verification).
  \item Check each concurrency boundary for cleanliness.
\end{enumerate}
No global model of the event-loop schedule is needed.
\end{corollary}

\begin{corollary}[Pure-coroutine race-freedom]\label{cor:pure}
If every coroutine of $P$ has effect class $\Epure$ (no shared-state
writes), then $P$ is trivially race-free: all concurrency boundaries are
clean by \cref{def:clean-cb}(i).
\end{corollary}

\begin{corollary}[Cancellation safety]\label{cor:cancel}
If every \texttt{await} point in $P$ is guarded by a \texttt{try}/\texttt{finally}
block that restores all shared variables to their pre-\texttt{await} values,
then the cancellation obstruction $\cancelObs$ is zero in Čech cohomology,
and $P$ is cancellation-safe.
\end{corollary}

\subsection{Limitations and decidability}

The theorem is stated for the \emph{static} analysis performed by
\texttt{AsyncAnalyzer} and \texttt{ConcurrencyBoundary}.  In general,
determining whether a concurrency boundary is clean requires reasoning about
dynamic aliasing, which is undecidable.  The implementation uses a sound
over-approximation: it conservatively marks a boundary as dirty whenever it
cannot prove cleanliness.  As a result, the theorem gives a sound but
incomplete verification: no false negatives (a clean program is never flagged
as racy), but false positives are possible.

% ======================================================================
\section{Experiments}\label{sec:experiments}
% ======================================================================

\subsection{Benchmark suite}

We evaluated the async effect analysis on a corpus of 120 Python files drawn
from popular asyncio-based libraries and internal Microsoft projects.  The
corpus was partitioned as follows:

\begin{table}[H]
  \centering
  \caption{Benchmark corpus statistics.}
  \label{tab:corpus}
  \begin{tabular}{lrrrr}
    \toprule
    Category & Files & Coroutines & Regions & Boundaries \\
    \midrule
    Pure async (no shared state) & 38 & 210 & 584 & 0 \\
    Shared-state (benign)        & 41 & 307 & 921 & 143 \\
    Shared-state (races)         & 21 & 198 & 612 & 89 \\
    Mixed IO                     & 20 & 174 & 503 & 61 \\
    \midrule
    Total                        & 120 & 889 & 2620 & 293 \\
    \bottomrule
  \end{tabular}
\end{table}

\subsection{Results}

\begin{table}[H]
  \centering
  \caption{Effect classification accuracy and boundary detection.}
  \label{tab:results}
  \begin{tabular}{lrrrr}
    \toprule
    Category & Effect acc.\ & CB clean & CB dirty & False pos. \\
    \midrule
    Pure async           & 100\% & n/a   & n/a  & 0 \\
    Shared-state (benign) & 97.3\% & 138  & 5    & 5 \\
    Shared-state (races) & 96.2\% & 31   & 58   & 3 \\
    Mixed IO             & 95.8\% & 55   & 6    & 6 \\
    \midrule
    Overall              & 97.1\% & 224  & 69   & 14 \\
    \bottomrule
  \end{tabular}
\end{table}

\paragraph{Effect classification.}
The \texttt{AsyncAnalyzer} achieves 97.1\% accuracy in assigning the correct
effect class to each async region.  The 2.9\% misclassifications arise
primarily from dynamic dispatch (effect class of a called function depends on
the runtime type of an argument) and reflection (\texttt{getattr} on a shared
object).

\paragraph{Boundary detection.}
Of the 293 concurrency boundaries, \texttt{ConcurrencyBoundary} correctly
identifies 69 dirty boundaries (all known races are flagged).  There are 14
false positives, all arising from over-approximation of the alias analysis:
the analyzer conservatively treats two variables as aliased when they are in
fact distinct objects.

\paragraph{Performance.}
Analysis time is sub-linear in the number of regions: on the full corpus of
2620 regions, the analysis completes in 1.7\,s on a standard laptop.  The
dominant cost is scope analysis (ScopeAnalyzer traversal), which accounts for
63\% of the total time.

\subsection{Case study: asyncio web crawler}

We applied the analysis to a 400-line asyncio web crawler with 18 coroutines
and 62 async regions.  The \texttt{AsyncAnalyzer} identified 4 dirty
boundaries, all involving a shared \texttt{visited\_urls} set.  Inspection
confirmed that two of these were genuine logical races; the other two were
false positives from URL-object aliasing.  After the developer added
\texttt{asyncio.Lock} guards, all four boundaries became clean and the
analysis passed.

\subsection{Comparison with existing tools}

\begin{table}[H]
  \centering
  \caption{Comparison with related async analysis tools.}
  \label{tab:comparison}
  \begin{tabular}{lcccc}
    \toprule
    Tool & Formal guarantee & Effect classes & Cancellation & Speed \\
    \midrule
    \JuGeo{} (this work) & Race-free theorem & 4           & Yes & $<2$\,s \\
    pyright             & Type checking only & 0           & No  & $<1$\,s \\
    mypy + await-lints  & Type checking only & 0           & No  & $<5$\,s \\
    asyncio-lock-checker & None             & 1 (stateful) & No  & $<1$\,s \\
    \bottomrule
  \end{tabular}
\end{table}

% ======================================================================
\section{Conclusion}\label{sec:conclusion}
% ======================================================================

We have presented a sheaf-theoretic framework for verifying Python's
\texttt{async}/\texttt{await} programs.  The key insight is that async
regions are naturally the coordinates of a Grothendieck site, and
\texttt{await} points are restriction morphisms.  This reformulation
transforms a global concurrency problem---is this program race-free?---into a
collection of local verification problems plus a boundary-cleanliness check.
The resulting modular proof principle (the Async Safety Theorem,
\cref{thm:async-safety}) reduces global verification to:
\begin{enumerate}
  \item Local effect annotation of each async region (done automatically by
        \texttt{AsyncAnalyzer}).
  \item A boundary scan by \texttt{ConcurrencyBoundary} that checks each
        pair of concurrently-reachable regions for conflicting state access.
\end{enumerate}

\noindent
The approach handles all four effect classes (pure, exception, stateful, IO),
tracks lexical scope and state through \texttt{ScopeAnalyzer} and
\texttt{StateTracker}, and models cancellation as a Čech 1-cocycle
obstruction.  Experiments on a 120-file corpus achieve 97.1\% effect
classification accuracy and detect all known races with 14 false positives.

\paragraph{Future work.}
Several directions remain open:
\begin{itemize}
  \item \textbf{Higher-order async.}  Coroutines that accept other coroutines
        as arguments require extending the site with higher-order
        morphisms---a direction related to the operadic composition of
        Paper~14.
  \item \textbf{Structured concurrency.}  Python 3.11+ \texttt{TaskGroup}
        introduces structured concurrency semantics; its lifetimes should be
        modeled as sheaves on an interval-order site.
  \item \textbf{Alias analysis.}  Reducing false positives requires a finer
        alias analysis; incorporating points-to information from the
        \texttt{StateTracker} is a natural next step.
  \item \textbf{Integration with Paper~3.}  Concurrency boundaries are
        obstructions in the Čech complex; future work should make this
        connection precise and use the descent machinery from Paper~3 to
        give completeness results.
\end{itemize}

\bigskip
\noindent\textbf{Acknowledgements.}
The Lean formalization was completed with assistance from the \JuGeo{}
proof assistant.  The \texttt{AsyncAnalyzer} and \texttt{ConcurrencyBoundary}
implementations were developed in collaboration with the JuGeo Python runtime
team.

\bibliographystyle{plain}
\begin{thebibliography}{99}

\bibitem{pep492}
Y.\,Selivanov.
\newblock {PEP 492 -- Coroutines with async and await syntax}.
\newblock Python Enhancement Proposal 492, 2015.

\bibitem{pep525}
Y.\,Selivanov.
\newblock {PEP 525 -- Asynchronous generators}.
\newblock Python Enhancement Proposal 525, 2016.

\bibitem{pep567}
Y.\,Selivanov.
\newblock {PEP 567 -- Context variables}.
\newblock Python Enhancement Proposal 567, 2018.

\bibitem{pep654}
I.\,Levkivskyi, C.\,Langa, G.\,van Rossum.
\newblock {PEP 654 -- Exception groups and except*}.
\newblock Python Enhancement Proposal 654, 2021.

\bibitem{sga4}
M.\,Artin, A.\,Grothendieck, J.-L.\,Verdier.
\newblock \textit{Th\'{e}orie des Topos et Cohomologie \'{E}tale des
  Sch\'{e}mas (SGA~4)}.
\newblock Lecture Notes in Mathematics 269, 270, 305.
\newblock Springer, 1972--1973.

\bibitem{maclane-moerdijk}
S.\,Mac Lane, I.\,Moerdijk.
\newblock \textit{Sheaves in Geometry and Logic}.
\newblock Springer, 1994.

\bibitem{johnstone}
P.\,T.\,Johnstone.
\newblock \textit{Sketches of an Elephant: A Topos Theory Compendium}.
\newblock Oxford University Press, 2002.

\bibitem{hyttinen}
T.\,Hyttinen, G.\,Oikkonen, J.\,V\"{a}\"{a}n\"{a}nen.
\newblock Sheaves and the concurrency problem.
\newblock \textit{Theoretical Computer Science}, 373(1--2):143--160, 2007.

\bibitem{jugeo-paper1}
H.\,Young.
\newblock Semantic sites for software verification.
\newblock Judgment Geometry Series, Paper~1, 2024.

\bibitem{jugeo-paper3}
H.\,Young.
\newblock Descent obstructions and sheaf cohomology.
\newblock Judgment Geometry Series, Paper~3, 2024.

\bibitem{jugeo-paper7}
H.\,Young.
\newblock Python effects as judgment sections.
\newblock Judgment Geometry Series, Paper~7, 2024.

\bibitem{jugeo-paper14}
H.\,Young.
\newblock Operadic composition of verification strategies.
\newblock Judgment Geometry Series, Paper~14, 2024.

\bibitem{trio}
N.\,J.\,Smith.
\newblock \textit{Notes on structured concurrency}.
\newblock \url{https://vorpus.org/blog/notes-on-structured-concurrency-or-go-statement-considered-harmful/}, 2018.

\bibitem{asyncio}
G.\,van Rossum.
\newblock \textit{asyncio --- Asynchronous I/O}, Python standard library
  documentation.
\newblock \url{https://docs.python.org/3/library/asyncio.html}, 2013.

\end{thebibliography}

\end{document}
